% ============================================================================
% chapter-template.tex  --  Scaffold for ONE chapter.
% Copy this into chapters/<name>.tex, then add  \include{chapters/<name>}
% to thesis.tex in reading order.  Do NOT add a preamble or \documentclass:
% this file is meant to be \include'd by the master thesis.tex.
%
% Replace "Chapter Title" and the label "ch:template" below.
% ============================================================================

\chapter{Chapter Title}
\label{ch:template}

This chapter introduces the topic. Reference other parts of the thesis with
cleveref so the words ("Chapter", "Figure", "Table") are added automatically,
e.g. background material appears in \cref{ch:template}.

% ---------------------------------------------------------------------------
\section{A Section}
\label{sec:template-section}

Write content here. Cite shared references by their key from
\texttt{references.bib}, for example \textcite{example2020key} or \cite{example2020key}.

% ---- A labelled figure ----------------------------------------------------
% Place the image file in figures/ . The extension can be omitted.
\begin{figure}[ht]
  \centering
  % \includegraphics[width=0.6\textwidth]{example-image}
  \fbox{\rule{0pt}{3cm}\rule{6cm}{0pt}}% placeholder box; swap for \includegraphics
  \caption{An example figure. Replace the placeholder box with \texttt{\textbackslash includegraphics}.}
  \label{fig:template-figure}
\end{figure}

The pipeline is shown in \cref{fig:template-figure}.

% ---- A labelled table -----------------------------------------------------
\begin{table}[ht]
  \centering
  \begin{tabular}{lr}
    \toprule
    Item & Value \\
    \midrule
    Alpha & 1.0 \\
    Beta  & 2.5 \\
    \bottomrule
  \end{tabular}
  \caption{An example table.}
  \label{tab:template-table}
\end{table}

Results are summarised in \cref{tab:template-table}.

% ---- A labelled equation --------------------------------------------------
\begin{equation}
  \label{eq:template-loss}
  \mathcal{L} = \frac{1}{n} \sum_{i=1}^{n} (y_i - \hat{y}_i)^2
\end{equation}

The objective is defined in \cref{eq:template-loss}.

% ---------------------------------------------------------------------------
\section{Another Section}
\label{sec:template-section-two}

Continue the chapter here. Cross-reference an earlier section with
\cref{sec:template-section}.
